\chapter{Introduction}

\section{Motivation}

The use of prosthetic devices in animals has increased in recent years, driven by advances in veterinary medicine and biomedical engineering. These devices aim to restore mobility and improve the quality of life of animals that have suffered limb loss or severe injury.

Despite these advances, the evaluation of prosthesis effectiveness and adaptation remains largely based on qualitative observation. Veterinarians typically rely on visual assessment of gait and general behavior, which introduces subjectivity and limits the consistency and reproducibility of the evaluation process.

This lack of objective and quantitative assessment methods represents a significant limitation, particularly in the context of long-term monitoring and optimization of prosthetic devices. The ability to measure locomotion using reliable and repeatable metrics is therefore essential to improve clinical decision-making and prosthesis design.

Furthermore, the development of quantitative evaluation methods in veterinary prosthetics may also contribute to advancements in human prosthetics, where similar challenges regarding objective performance assessment are present.

\section{Topic Overview}

This work proposes the development of an instrumented prosthesis for canine locomotion monitoring. The system integrates inertial and force sensors to acquire biomechanical data during movement, enabling the characterization of locomotor behavior through measurable parameters.

The acquired data is processed and analyzed to extract relevant features related to movement consistency, force distribution, and temporal patterns. These features provide a basis for evaluating the adaptation of the animal to the prosthesis.

The system is designed to operate in real-world conditions, allowing data acquisition outside controlled laboratory environments and supporting continuous monitoring over extended periods.

\section{Objectives and Deliverables}

The main objective of this work is to develop a system capable of providing objective and quantitative assessment of prosthesis adaptation in animals.

This objective can be divided into the following specific goals:

\begin{itemize}
    \item Development of an instrumented prosthesis integrating inertial and force sensors;
    \item Implementation of a data acquisition system for continuous monitoring of locomotion;
    \item Definition of biomechanical metrics to characterize locomotor behavior;
    \item Exploration of data-driven methods, including machine learning, to classify adaptation states;
    \item Validation of the proposed approach through experimental data.
\end{itemize}

The expected deliverables include a functional prototype, a dataset of acquired signals, and a methodology for analyzing locomotion and assessing prosthesis adaptation.

\section{Thesis Outline}

This document is organized as follows. Chapter 2 presents the background concepts related to canine locomotion, biomechanical analysis, and sensing technologies. Chapter 3 reviews the current state of the problem and discusses the limitations of existing approaches. Chapter 4 describes the proposed system and its main components.

Chapter 5 details the methodology used for data acquisition and analysis, including the definition of relevant metrics and the formulation of the classification problem. Chapter 6 presents the preliminary work and initial results obtained. Chapter 7 outlines the work plan and future development steps. Finally, Chapter 8 concludes the document and discusses potential future work.